\section*{Step 314 Updated Theorem Status}

\paragraph{Derivative-ratio phase formula.}
For
\[
T_{k,j}={k\choose j}\zeta^{(j)}(\rho_1)M(G_\star)^{(k-j)}(\rho_1),
\]
the binomial factor is positive real, hence
\[
\arg T_{k,j}
=
\arg\zeta^{(j)}(\rho_1)
+
\arg M(G_\star)^{(k-j)}(\rho_1).
\]
The forward phase increment is
\[
\arg\frac{T_{k,j+1}}{T_{k,j}}
=
\arg\frac{\zeta^{(j+1)}(\rho_1)}{\zeta^{(j)}(\rho_1)}
+
\arg\frac{M(G_\star)^{(k-j-1)}(\rho_1)}
        {M(G_\star)^{(k-j)}(\rho_1)}.
\]
The continuous-index approximation requested in Step 314 is
\[
\alpha(k)\approx
\arg\frac{\zeta^{(k/2+1)}(\rho_1)}{\zeta^{(k/2)}(\rho_1)}
-
\arg\frac{M(G_\star)^{(k/2+1)}(\rho_1)}
        {M(G_\star)^{(k/2)}(\rho_1)}.
\]

\paragraph{Numerical status.}
This formula matches the empirical central phase slopes to a few percent
on \(k=10,20,30,50\), and the symmetric exact adjacent-ratio formula
matches the local Step 313 phase derivative especially well for
\(k\ge20\).

\paragraph{Rigorous status.}
The theorem is not upgraded.  The missing theorem is a uniform bound or
asymptotic for
\[
\arg\frac{\zeta^{(j+1)}(\rho_1)}{\zeta^{(j)}(\rho_1)}
\]
near \(j\sim k/2\), plus the corresponding uniform control of the
Mellin-derivative ratio.  Existing literature located in this audit
addresses moments and ratios-conjecture averages of zeta derivatives at
zeros; it does not supply the fixed-zero, high-order phase-ratio bound
needed here.

